% =============================================================================
%  Slides -- Projeto Final PCR / Co-Projeto HW-SW -- UnB 2026/1
%  Apresentacao de 20 min.
%
%  Compilacao: pdflatex slides_pcr.tex  (2x)
%  Requer os .dat em docs/figs/ (gerados por golden/gen_figs.py)
%
%  ==> Slides marcados com [VIVADO] precisam dos prints/numeros da sintese.
% =============================================================================
\documentclass[aspectratio=169, 11pt]{beamer}

\usepackage[utf8]{inputenc}
\usepackage[T1]{fontenc}
\usepackage[brazilian]{babel}
\usepackage{booktabs}
\usepackage{tikz}
\usepackage{pgfplots}
\usepackage{siunitx}
\pgfplotsset{compat=1.17}
\usetikzlibrary{arrows.meta, positioning, calc}
\sisetup{output-decimal-marker={,}}

\usetheme{metropolis}
\metroset{progressbar=frametitle, numbering=fraction}
\definecolor{unb}{RGB}{0,70,130}
\setbeamercolor{frametitle}{bg=unb}
\setbeamercolor{alerted text}{fg=red!75!black}

\title{Acelerador em FPGA do Supervisor LiDAR}
\subtitle{Co-Projeto HW--SW sobre Zynq-7000 \\ \small Cadeira de rodas semi-assistida}
\author{Vítor Gonçalves de Andrade Silva \and Thiago Henrique Oliveira Souza}
\institute{Engenharia Eletrônica --- FCTE/UnB}
\date{2026/1}

\begin{document}

\maketitle

% =============================================================================
\begin{frame}{O sistema}
\begin{columns}[T]
\begin{column}{0.55\textwidth}
Cadeira de rodas motorizada \textbf{semi-assistida}.

\vspace{4mm}
\alert{Princípio de segurança:}
\begin{center}
\large o sistema pode \textbf{FREAR} ou \textbf{DESVIAR},\\
mas \textbf{nunca ACELERAR}.
\end{center}
\vspace{2mm}
A iniciativa é sempre do usuário.

\vspace{4mm}
\footnotesize
\begin{itemize}
\item Raspberry Pi 5 + ROS\,2 (percepção)
\item ESP32-S3 (tempo real: joystick, motores, encoders)
\item RPLIDAR C1 a \SI{10}{\hertz}, \num{720} feixes
\end{itemize}
\end{column}
\begin{column}{0.45\textwidth}
\begin{block}{O supervisor}
A cada varredura, decide se o avanço comandado deve ser \textbf{reduzido} ou
\textbf{interrompido}.
\end{block}
\vspace{2mm}
\begin{alertblock}{Por que acelerar isto?}
É o nó mais pesado \textbf{e} o único cuja latência tem \textbf{consequência
física}: ela vira \textbf{centímetros} rumo ao obstáculo.
\end{alertblock}
\end{column}
\end{columns}
\end{frame}

% =============================================================================
\begin{frame}{O algoritmo: teste de corredor}
\begin{columns}[T]
\begin{column}{0.42\textwidth}
\footnotesize
Projeta cada feixe no referencial do veículo:
\begin{align*}
X &= x_L + r\cos(a + \theta_L)\\
Y &= y_L + r\sin(a + \theta_L)
\end{align*}

É obstáculo se cair no \textbf{corredor que a cadeira vai ocupar}:
\[ |Y| \le W/2 \quad\text{e}\quad X > X_f \]

\vspace{3mm}
Para cada um dos \alert{$K=19$} candidatos de desvio ($-45^\circ$ a $+45^\circ$),
gira-se o ponto e repete-se o teste.
\end{column}
\begin{column}{0.58\textwidth}
\begin{center}
\begin{tikzpicture}[scale=0.85]
\begin{axis}[width=7.2cm, height=5cm,
  xlabel={$X$ (m)}, ylabel={$Y$ (m)},
  xmin=0, xmax=2.6, ymin=-1.0, ymax=1.3,
  tick label style={font=\tiny}, label style={font=\tiny},
  grid=major, grid style={gray!20}]
  \addplot[only marks, mark=*, mark size=0.4pt, gray!50]
    table[x=x, y=y] {figs/scan_out.dat};
  \addplot[only marks, mark=*, mark size=1pt, red]
    table[x=x, y=y] {figs/scan_in.dat};
  \draw[blue, very thick] (axis cs:0.687,-0.305) rectangle (axis cs:2.6,0.305);
  \draw[fill=black] (axis cs:0.667,0.336) circle (2pt);
  \node[font=\tiny, anchor=south] at (axis cs:0.667,0.38) {LiDAR};
\end{axis}
\end{tikzpicture}
\end{center}
\centering\tiny Varredura \textbf{real} da cadeira. Vermelho: no corredor.
\end{column}
\end{columns}

\vspace{2mm}
\centering
\alert{\textbf{$N \times K = 720 \times 19 = 9.918$ avaliações por varredura}}
\end{frame}

% =============================================================================
\begin{frame}{Por que este bloco: um defeito real de segurança}
\begin{columns}[T]
\begin{column}{0.5\textwidth}
\footnotesize
O supervisor \textbf{antigo} media a folga com um \textbf{cone} ancorado no
sensor.

\vspace{2mm}
Só valeria se o LiDAR estivesse \textbf{no eixo} da cadeira.

\vspace{2mm}
Ele está \alert{\SI{33.6}{\centi\metre} fora do eixo}.

\vspace{3mm}
\begin{alertblock}{Consequência}
A metade \textbf{direita} da cadeira ficava \textbf{cega} --- e a cobertura
\textbf{piora} quanto mais \textbf{perto} está o obstáculo.

\vspace{1mm}
A cadeira \textbf{atropelava} obstáculos à direita.
\end{alertblock}
\end{column}
\begin{column}{0.5\textwidth}
\begin{center}
\begin{tikzpicture}
\begin{axis}[width=6.2cm, height=4.6cm,
  xlabel={distância do obstáculo (m)}, ylabel={cobertura (\%)},
  ymin=0, ymax=105, xmin=0.3, xmax=1.5,
  tick label style={font=\tiny}, label style={font=\tiny},
  legend style={font=\tiny, at={(0.98,0.03)}, anchor=south east},
  grid=major, grid style={gray!20}]
  \addplot[very thick, red] table[x=d, y=cone30] {figs/cobertura.dat};
  \addlegendentry{cone $\pm30^\circ$}
  \addplot[very thick, orange, dashed] table[x=d, y=cone20] {figs/cobertura.dat};
  \addlegendentry{cone $\pm20^\circ$}
  \addplot[very thick, blue, densely dotted] coordinates {(0.3,100) (1.5,100)};
  \addlegendentry{corredor}
\end{axis}
\end{tikzpicture}
\end{center}
\centering\tiny A \SI{0.60}{\metre} (zona de parada) o cone vê só \SI{52}{\percent}.
\end{column}
\end{columns}
\end{frame}

% =============================================================================
\begin{frame}{Particionamento HW/SW}
\begin{center}
\begin{tabular}{@{}llc@{}}
\toprule
\textbf{Etapa} & \textbf{Custo} & \textbf{Onde} \\
\midrule
Projeção + teste de corredor & $O(N\cdot K) = \mathbf{9.918}$ & \alert{\textbf{FPGA}} \\
Função custo + \emph{argmin} & $O(K) = 19$ & ARM \\
\bottomrule
\end{tabular}
\end{center}

\vspace{5mm}
\begin{block}{Por que \emph{não} levar tudo para hardware?}
A função custo são \textbf{19 comparações} --- o ganho seria irrisório.

Em compensação, é nela que moram os \textbf{pesos ajustados em campo}
($w_{\text{obs}}$, $w_{\text{dev}}$, $K_{\text{assist}}$), retocados a cada ensaio.

\vspace{2mm}
\centering\alert{Congelá-los em RTL trocaria flexibilidade por nada.}
\end{block}

\vspace{2mm}
\centering\footnotesize
É exatamente o \emph{trade-off} que o co-projeto existe para arbitrar.
\end{frame}

% =============================================================================
\begin{frame}{Arquitetura de hardware}
\begin{center}
\begin{tikzpicture}[font=\small, node distance=8mm,
  blk/.style={draw, rounded corners=2pt, minimum height=11mm, minimum width=20mm,
              inner sep=4pt, align=center},
  arr/.style={-Latex, thick}]
  \node[blk, fill=gray!10] (in) {AXI4-Stream\\\scriptsize $(r_i,\theta_i)$};
  \node[blk, right=of in, fill=blue!12] (cor) {CORDIC\\\scriptsize 16 estágios\\\alert{\scriptsize 0 DSP}};
  \node[blk, right=of cor, fill=gray!10] (prj) {projeção\\\scriptsize $X_i,\,Y_i$};
  \node[blk, right=of prj, fill=orange!18, minimum height=20mm] (lanes)
       {\textbf{19 pistas}\\\textbf{paralelas}\\\alert{\scriptsize 76 DSP}};
  \node[blk, right=of lanes, fill=gray!10] (out) {AXI4-Lite\\\scriptsize 19 folgas};
  \draw[arr] (in) -- (cor);
  \draw[arr] (cor) -- node[above, font=\tiny] {$\cos,\sin$} (prj);
  \draw[arr] (prj) -- (lanes);
  \draw[arr] (lanes) -- (out);
\end{tikzpicture}
\end{center}

\vspace{2mm}
\begin{columns}[T]
\begin{column}{0.5\textwidth}
\footnotesize
\textbf{CORDIC compartilhado} ($O(N)$)
\begin{itemize}\footnotesize
\item Só somadores e \emph{shifts} $\Rightarrow$ \alert{zero DSP}
\item \emph{Por que não LUT de seno?} O ângulo depende da \textbf{calibração},
reescrita em tempo de execução via AXI4-Lite.
\end{itemize}
\end{column}
\begin{column}{0.5\textwidth}
\footnotesize
\textbf{19 pistas replicadas} ($O(N\cdot K)$)
\begin{itemize}\footnotesize
\item $\cos\delta_k,\sin\delta_k$ são \textbf{constantes de síntese}
\item $\Rightarrow$ 4 multiplicadores fixos por pista
\item Avaliam \textbf{no mesmo ciclo} o que o software \textbf{itera}
\end{itemize}
\end{column}
\end{columns}

\vspace{2mm}
\centering\alert{Vazão: \textbf{1 feixe/ciclo}, sem bolhas.}
\end{frame}

% =============================================================================
\begin{frame}{Verificação: equivalência bit a bit}
\begin{columns}[T]
\begin{column}{0.52\textwidth}
\footnotesize
\begin{enumerate}
\item Modelo \textbf{ponto fixo} (Python) emula o RTL \emph{exatamente}
\item Modelo \textbf{float} = referência física
\item O modelo gera \textbf{estímulo} e \textbf{esperado}
\item Testbench VHDL lê os arquivos e compara as 19 folgas
\end{enumerate}

\vspace{3mm}
\begin{block}{\footnotesize Detalhe que evita um erro invisível}
\scriptsize As constantes do RTL (\texttt{cordic\_pkg.vhd}) são \textbf{geradas}
pelo mesmo modelo. RTL e referência casam \textbf{por construção} --- não por
transcrição manual.
\end{block}
\end{column}
\begin{column}{0.48\textwidth}
\begin{exampleblock}{Resultado (GHDL)}
\scriptsize
\texttt{varredura 0: EQUIVALENCIA OK}\\
\texttt{varredura 1: EQUIVALENCIA OK}\\
\texttt{varredura 2: EQUIVALENCIA OK}
\vspace{1mm}

\normalsize\centering
\textbf{Zero divergência}
\end{exampleblock}

\vspace{2mm}
\footnotesize
O estímulo \textbf{não é sintético}: são três varreduras \alert{reais} da cadeira,
com um balde a \SI{21}{\centi\metre} à direita --- o antigo ponto cego.
\end{column}
\end{columns}
\end{frame}

% =============================================================================
\begin{frame}{Precisão: o achado do trabalho}
\begin{center}
\alert{\Large O primeiro modelo em ponto fixo errava \textbf{\SI{402}{\milli\metre}}.}
\end{center}

\vspace{3mm}
\begin{columns}[T]
\begin{column}{0.52\textwidth}
\footnotesize
Não era arredondamento. Era \textbf{estouro de faixa}:
\[ \theta = a_{\min} + i\,\Delta a + \theta_L \;\to\; \alert{\SI{-5.85}{\radian}} \]
\vspace{-3mm}
\[ \text{Q3.13 comporta apenas } \pm\SI{4.0}{\radian} \]

\vspace{2mm}
O ângulo \textbf{saturava} $\Rightarrow$ o feixe era projetado numa direção
completamente errada.

\vspace{2mm}
\textbf{Correção:} envolver $\theta$ em $[-\pi,\pi)$.
\end{column}
\begin{column}{0.48\textwidth}
\centering
\begin{tabular}{@{}lrr@{}}
\toprule
& \scriptsize médio & \scriptsize máximo \\
\midrule
\scriptsize $\theta$ saturado & \scriptsize \SI{17.89}{\milli\metre} & \alert{\scriptsize \SI{402}{\milli\metre}} \\
\scriptsize $\theta$ envolvido & \scriptsize \textbf{\SI{0.99}{\milli\metre}} & \scriptsize \textbf{\SI{5.27}{\milli\metre}} \\
\bottomrule
\end{tabular}

\vspace{4mm}
\footnotesize
LSB de Q6.10 $=$ \SI{0.98}{\milli\metre}

$\Rightarrow$ erro final $\approx$ \textbf{1 LSB}

$\Rightarrow$ \SI{0.9}{\percent} da zona de parada
\end{column}
\end{columns}

\vspace{4mm}
\begin{alertblock}{Lição}
\centering
Em ponto fixo, a \textbf{FAIXA} costuma ser mais crítica que a \textbf{RESOLUÇÃO}.\\
\footnotesize Perdemos \SI{402}{\milli\metre} por 3 bits de parte inteira --- não por falta de bits fracionários.
\end{alertblock}
\end{frame}

% =============================================================================
\begin{frame}{Saída do HW vs saída do SW}
\begin{columns}[T]
\begin{column}{0.52\textwidth}
\begin{center}
\begin{tikzpicture}
\begin{axis}[width=6.4cm, height=4.2cm,
  xlabel={$\delta_k$ (graus)}, ylabel={folga (m)},
  xmin=-48, xmax=48, xtick={-45,-15,15,45},
  legend style={font=\tiny, at={(0.02,0.98)}, anchor=north west},
  grid=major, grid style={gray!20},
  tick label style={font=\tiny}, label style={font=\tiny}]
  \addplot[thick, blue, mark=o, mark size=2pt]
    table[x=delta, y=sw] {figs/sw_hw_scan0.dat};
  \addlegendentry{SW (float)}
  \addplot[red, only marks, mark=x, mark size=3.5pt]
    table[x=delta, y=hw] {figs/sw_hw_scan0.dat};
  \addlegendentry{HW (RTL simulado)}
\end{axis}
\end{tikzpicture}
\end{center}
\centering\tiny As curvas \textbf{se sobrepõem} --- é esse o resultado.
\end{column}
\begin{column}{0.48\textwidth}
\begin{center}
\begin{tikzpicture}
\begin{axis}[width=5.8cm, height=4.2cm,
  xlabel={$\delta_k$ (graus)}, ylabel={erro HW$-$SW (mm)},
  xmin=-48, xmax=48, xtick={-45,-15,15,45},
  legend style={font=\tiny, at={(0.98,0.02)}, anchor=south east},
  grid=major, grid style={gray!20},
  tick label style={font=\tiny}, label style={font=\tiny}]
  \draw[gray, dashed] (axis cs:-48,0.98) -- (axis cs:48,0.98);
  \draw[gray, dashed] (axis cs:-48,-0.98) -- (axis cs:48,-0.98);
  \addplot[only marks, mark=*, mark size=1.4pt, blue]
    table[x=delta, y=err_mm] {figs/err_scan0.dat};
  \addplot[only marks, mark=square*, mark size=1.4pt, red]
    table[x=delta, y=err_mm] {figs/err_scan1.dat};
  \addplot[only marks, mark=triangle*, mark size=1.8pt, orange]
    table[x=delta, y=err_mm] {figs/err_scan2.dat};
\end{axis}
\end{tikzpicture}
\end{center}
\centering\tiny Tracejado: $\pm 1$ LSB. 3 varreduras, 57 comparações.
\end{column}
\end{columns}

\vspace{1mm}
\centering\footnotesize
Erro médio \textbf{\SI{0.99}{\milli\metre}} $\;\vert\;$ máximo \textbf{\SI{5.27}{\milli\metre}}
$\;\vert\;$ \alert{$\approx$ 1 LSB} $\;\vert\;$ \SI{0.9}{\percent} da zona de parada

\vspace{1mm}
\tiny O ``HW'' é a saída \textbf{real} da simulação VHDL (\texttt{rtl\_out}), não o modelo.
\end{frame}

% =============================================================================
\begin{frame}{Desempenho}
\begin{center}
\begin{tabular}{@{}lrrr@{}}
\toprule
\textbf{Implementação} & \textbf{Tempo/varredura} & \multicolumn{2}{c}{\textbf{Speedup}} \\
 & & \scriptsize vs escalar & \scriptsize vs NumPy \\
\midrule
Software escalar (Pi 5) & \SI{3.54}{\milli\second} & --- & --- \\
Software NumPy (Pi 5) & \SI{0.488}{\milli\second} & $7{,}3\times$ & --- \\
\alert{\textbf{HW \SI{100}{\mega\hertz}}} \scriptsize(751 ciclos) & \alert{\textbf{\SI{7.51}{\micro\second}}} & \alert{$\mathbf{472\times}$} & \alert{$\mathbf{65\times}$} \\
HW \SI{150}{\mega\hertz} & \SI{5.01}{\micro\second} & $707\times$ & $98\times$ \\
\bottomrule
\end{tabular}
\end{center}

\vspace{-1mm}
\begin{center}
\begin{tikzpicture}
\begin{semilogxaxis}[width=10.5cm, height=3.4cm,
  xlabel={tempo por varredura ($\mu$s --- escala log)},
  xmin=3, xmax=9000, xmajorgrids, grid style={gray!25},
  symbolic y coords={HW150, HW100, NumPy, escalar},
  ytick=data, y dir=reverse,
  yticklabels={HW \SI{150}{\mega\hertz}, HW \SI{100}{\mega\hertz}, SW NumPy, SW escalar},
  tick label style={font=\tiny}, label style={font=\tiny},
  nodes near coords, nodes near coords style={font=\tiny, anchor=west},
  point meta=explicit symbolic,
  xbar, bar width=7pt, enlarge y limits=0.28]
  \addplot[fill=red!30, draw=red!70!black] coordinates {
    (3542, escalar) [\SI{3.54}{\milli\second}]
    (488,  NumPy)   [\SI{0.488}{\milli\second}]};
  \addplot[fill=blue!30, draw=blue!70!black] coordinates {
    (7.51, HW100)   [\SI{7.51}{\micro\second}]
    (5.01, HW150)   [\SI{5.01}{\micro\second}]};
\end{semilogxaxis}
\end{tikzpicture}
\end{center}
\centering\tiny
Escala \textbf{logarítmica} --- quase três ordens de grandeza. \emph{Baseline}
medido na Pi 5; ciclos de HW \textbf{medidos na simulação}.

\vspace{2mm}
\begin{alertblock}{Análise crítica: o ganho de tempo \emph{não era necessário}}
\footnotesize
O LiDAR entrega uma varredura a cada \SI{100}{\milli\second}. Mesmo o software
vetorizado usa só \textbf{\SI{0.5}{\percent}} desse orçamento --- \textbf{não
estávamos perdendo prazos}.
\end{alertblock}
\end{frame}

% =============================================================================
\begin{frame}{Então qual é o valor real?}
\begin{enumerate}
\item \textbf{Alívio de CPU} \\
      \footnotesize Libera \SI{8.3}{\percent} numa Pi que também roda SLAM e EKF ---
      e que já \textbf{congelou} por exaustão de memória.
      \vspace{3mm}
\item \textbf{Latência determinística} \alert{$\leftarrow$ o mais importante} \\
      \footnotesize Um nó Python está sujeito ao escalonador de um Linux
      não-preemptivo. Num laço de \textbf{segurança}, determinismo vale mais que
      média baixa.
      \vspace{3mm}
\item \textbf{Orçamento para crescer} \\
      \footnotesize Dá para aumentar $K$ ou a densidade angular sem custo de tempo:
      o gargalo passa a ser o \textbf{sensor}.
\end{enumerate}

\vspace{4mm}
\centering
\footnotesize Reportamos isso explicitamente porque a leitura oposta ---
``aceleramos, logo melhorou'' --- seria \textbf{desonesta com os dados}.
\end{frame}

% =============================================================================
\begin{frame}{Síntese: recursos e timing \hfill \scriptsize Vivado 2022.2, xc7z020}
\begin{columns}[T]
\begin{column}{0.46\textwidth}
\centering
\begin{tabular}{@{}lrrr@{}}
\toprule
 & \scriptsize prev. & \scriptsize obtido & \scriptsize \% \\
\midrule
LUT & --- & 2.447 & 4,6 \\
FF & --- & 1.792 & 1,7 \\
\textbf{DSP48} & 78 & \alert{\textbf{54}} & 24,5 \\
BRAM & 0 & \textbf{0} & 0 \\
\bottomrule
\end{tabular}

\vspace{3mm}
\begin{exampleblock}{\footnotesize Timing}
\centering\footnotesize
WNS $=$ \textbf{\SI{+0.872}{\nano\second}}\\
$\Rightarrow$ \alert{fecha a \SI{100}{\mega\hertz}}\\[1mm]
$F_{\max} = $ \textbf{\SI{109.6}{\mega\hertz}}
\end{exampleblock}
\centering\tiny Potência: ver \texttt{power.rpt} \hfill 0 erros, 0 \emph{critical warnings}
\end{column}
\begin{column}{0.54\textwidth}
\begin{alertblock}{Erramos a previsão de DSP --- para MAIS}
\footnotesize
Previmos \textbf{78} ($19\times4 + 2$). A ferramenta usou \textbf{54} --- \SI{31}{\percent} a menos.
\end{alertblock}

\footnotesize
\textbf{Por quê?} O DSP48E1 não é um multiplicador isolado: é um
\textbf{multiplicador-acumulador} com porta $C$ e cascata \texttt{PCIN}.

\vspace{1mm}
A síntese reconheceu que
\[ x^{(k)} = X\cos\delta_k + Y\sin\delta_k \]
é a forma que o bloco implementa \textbf{nativamente}, e fundiu produto $+$ soma
num \emph{único} DSP.

\vspace{2mm}
\alert{\footnotesize Contar operadores do RTL \textbf{superestima} o custo:
o mapeamento explora estrutura do dispositivo que não aparece no código.}
\end{column}
\end{columns}
\end{frame}

% =============================================================================
\begin{frame}{O que a síntese confirmou}
\begin{columns}[T]
\begin{column}{0.5\textwidth}
\begin{exampleblock}{CORDIC: \textbf{0 DSP}}
\footnotesize
Os \textbf{54} DSP estão \emph{todos} no \texttt{corridor\_core} (pistas $+$
projeção).

\vspace{1mm}
O CORDIC usa só somadores e deslocamentos cabeados --- exatamente o que se
pretendia ao escolhê-lo.
\end{exampleblock}

\begin{exampleblock}{BRAM: \textbf{0}}
\footnotesize
A arquitetura \emph{streaming} não precisa armazenar a varredura: cada feixe é
consumido \textbf{no ciclo em que chega}.
\end{exampleblock}
\end{column}
\begin{column}{0.5\textwidth}
\begin{block}{Folga para crescer}
\footnotesize
Sobram \textbf{\SI{75}{\percent}} dos DSP.

\vspace{2mm}
Dá para \textbf{triplicar} o número de candidatos $K$ sem sair do
\texttt{xc7z020} --- ou aumentar a precisão das pistas.

\vspace{2mm}
E o gargalo continuaria sendo o \textbf{sensor} (\SI{100}{\milli\second}), não a
computação.
\end{block}

\vspace{2mm}
\centering
\alert{\scriptsize [COLAR AQUI: print do \texttt{timing\_summary} \\ e do \texttt{power.rpt}]}
\end{column}
\end{columns}
\end{frame}

% =============================================================================
\begin{frame}{Block design: ARM + AXI-DMA + acelerador \hfill [VIVADO]}
\begin{center}
\begin{tikzpicture}[font=\scriptsize, node distance=10mm,
  blk/.style={draw, rounded corners=2pt, minimum height=13mm, minimum width=26mm,
              align=center, inner sep=4pt},
  arr/.style={-Latex, thick}]
  \node[blk, fill=green!12] (ps) {\textbf{Zynq PS}\\ARM Cortex-A9\\\scriptsize DDR: /scan};
  \node[blk, right=22mm of ps, fill=blue!12] (dma) {\textbf{AXI-DMA}};
  \node[blk, right=18mm of dma, fill=orange!18] (acc) {\textbf{lidar\_accel}\\\scriptsize CORDIC + 19 pistas};
  \draw[arr] (ps) -- node[above, font=\tiny] {AXI4-Lite} node[below, font=\tiny] {(GP0)} (dma);
  \draw[arr] (dma) -- node[above, font=\tiny] {AXI4-Stream} (acc);
  \draw[arr] (ps.south) -- ++(0,-7mm) -| node[pos=0.25, below, font=\tiny] {AXI4-Lite: calibração + 19 folgas} (acc.south);
  \draw[arr, red] (acc.north) -- ++(0,7mm) -| node[pos=0.25, above, font=\tiny] {IRQ} (ps.north);
\end{tikzpicture}
\end{center}

\vspace{3mm}
\footnotesize
\begin{itemize}
\item O ARM guarda a varredura na \textbf{DDR}; o \textbf{AXI-DMA} a transmite como
      \emph{stream} ao acelerador
\item Resultado (19 folgas) volta por \textbf{AXI4-Lite}; a \textbf{IRQ} avisa o fim
      --- sem espera ocupada
\item \emph{Stream} para o fluxo volumoso; mapeado em memória para os escalares.
\end{itemize}

\vspace{2mm}
\centering
\alert{\scriptsize [COLAR AQUI: print do block design do Vivado]}
\end{frame}

% =============================================================================
\begin{frame}{Conclusões}
\begin{block}{O que foi alcançado}
\footnotesize
\alert{$472\times$} vs software escalar, \alert{$65\times$} vs vetorizado, com
equivalência \textbf{bit a bit} verificada sobre \textbf{dados reais} da cadeira.
\end{block}

\begin{alertblock}{Análise crítica --- esperado vs alcançado}
\footnotesize
O \emph{speedup} veio, mas \textbf{o ganho de tempo não era necessário}: o software
já cabia no orçamento do sensor. O benefício legítimo é o \textbf{alívio de CPU} e,
sobretudo, o \textbf{determinismo} da latência num laço de segurança.
\end{alertblock}

\begin{exampleblock}{O resultado mais transferível}
\footnotesize
Não é o \emph{speedup}: é o \textbf{defeito de faixa em Q3.13}. Um ângulo de
\SI{-5.85}{\radian} saturando num formato de $\pm\SI{4}{\radian}$ produziu
\SI{402}{\milli\metre} de erro \textbf{silencioso} --- o RTL não acusa, a simulação
``roda''.

\vspace{1mm}
\centering\alert{Foi o \emph{golden model}, não a inspeção de ondas, que o encontrou.}
\end{exampleblock}
\end{frame}

% =============================================================================
\begin{frame}[standout]
Obrigado!

\vspace{5mm}
\small
\texttt{fpga/} --- RTL, golden model, testbench, scripts Vivado

\vspace{2mm}
\texttt{bash sim/run\_ghdl.sh} \\
\scriptsize reproduz a prova de equivalência em um comando
\end{frame}

% =============================================================================
\appendix
\begin{frame}{Backup: formato de ponto fixo}
\centering
\begin{tabular}{@{}llll@{}}
\toprule
\textbf{Grandeza} & \textbf{Formato} & \textbf{Faixa} & \textbf{LSB} \\
\midrule
distâncias, coords. & Q6.10 & $\pm\num{32}$\,m & \SI{0.98}{\milli\metre} \\
$\cos$, $\sin$ & Q1.14 & $\pm 2$ & \num{6.1e-5} \\
ângulos & Q3.13 & \alert{$\pm\SI{4.0}{\radian}$} & \SI{1.22e-4}{\radian} \\
\bottomrule
\end{tabular}

\vspace{6mm}
\footnotesize
A faixa de Q3.13 (\alert{$\pm\SI{4}{\radian}$}) foi a origem do defeito de
\SI{402}{\milli\metre}: $\theta$ chega a \SI{-5.85}{\radian}.
\end{frame}

\begin{frame}{Backup: por que CORDIC e não LUT?}
\footnotesize
Uma LUT de seno exigiria uma ROM indexada pelo \textbf{índice do feixe}.

\vspace{2mm}
Isso só funcionaria se \texttt{angle\_min}, o incremento e o \emph{yaw} fossem
\textbf{fixos em síntese}.

\vspace{2mm}
Mas eles são \alert{parâmetros de calibração} --- reescritos em tempo de execução
por AXI4-Lite sempre que o LiDAR é remontado.

\vspace{4mm}
\begin{block}{CORDIC}
\footnotesize
\begin{itemize}
\item Aceita \textbf{qualquer} ângulo em tempo de execução
\item Só somadores e deslocamentos cabeados $\Rightarrow$ \alert{zero DSP48}
\item 16 estágios pipelinados $\Rightarrow$ 1 resultado/ciclo
\item Redução de quadrante na entrada (o laço só converge para
      $|z| \le \SI{1.7434}{\radian}$)
\end{itemize}
\end{block}
\end{frame}

\end{document}
